% ════════════════════════════════════════════════════════════════════════════
% GUÍA — LOGARITMOS I: DEFINICIÓN Y CONEXIÓN CON POTENCIAS
% Semana 16 · 2026-S1 · Máximo + Martina · 2° Medio
% Generada a partir de ../../templates/guia-base.tex
% ════════════════════════════════════════════════════════════════════════════

\documentclass[11pt,a4paper]{article}
\usepackage[utf8]{inputenc}
\usepackage[T1]{fontenc}
\usepackage[spanish]{babel}
\usepackage{geometry}
\usepackage{amsmath,amssymb}
\usepackage{xcolor}
\usepackage{tcolorbox}
\usepackage{enumitem}
\usepackage{fancyhdr}
\usepackage{titlesec}
\usepackage{multirow}
\usepackage{array}
\usepackage{booktabs}
\usepackage{tikz}
\usepackage{hyperref}

\geometry{margin=2cm, top=2.5cm, bottom=2.5cm}

% ─── VARIABLES DE LA GUÍA ───
\newcommand{\guiaTitulo}{Logaritmos I}
\newcommand{\guiaSubtitulo}{Definición y conexión con potencias}
\newcommand{\guiaUnidad}{Unidad de Álgebra}
\newcommand{\guiaCurso}{2° Medio --- Matemáticas}
\newcommand{\guiaAlumno}{\underline{\hspace{6cm}}}
\newcommand{\guiaSemana}{Semana 16}
\newcommand{\guiaSemestre}{2026-S1}

% ─── Colores ───
\definecolor{azulprincipal}{HTML}{2980B9}
\definecolor{azuloscuro}{HTML}{1A5276}
\definecolor{verdequimica}{HTML}{1ABC9C}
\definecolor{verdeclaro}{HTML}{D5F5E3}
\definecolor{naranja}{HTML}{E67E22}
\definecolor{rojo}{HTML}{E74C3C}
\definecolor{amarillorecuerda}{HTML}{FFF3CD}
\definecolor{amarilloborde}{HTML}{FFC107}
\definecolor{grisfondo}{HTML}{F5F5F5}

% ─── Cajas ───
\tcbset{
  recuerda/.style={
    colback=amarillorecuerda, colframe=amarilloborde,
    fonttitle=\bfseries, title={\color{black} RECUERDA},
    boxrule=1pt, arc=2mm, left=5pt, right=5pt
  },
  propiedad/.style={
    colback=verdeclaro!50, colframe=verdequimica!80!black,
    fonttitle=\bfseries, boxrule=0.8pt, arc=2mm, left=5pt, right=5pt
  },
  formula/.style={
    colback=grisfondo, colframe=azulprincipal,
    boxrule=0.8pt, arc=2mm, left=5pt, right=5pt
  },
  definicion/.style={
    colback=blue!5, colframe=azulprincipal,
    fonttitle=\bfseries, boxrule=0.8pt, arc=2mm, left=5pt, right=5pt
  },
  ejemplo/.style={
    colback=orange!5, colframe=naranja,
    fonttitle=\bfseries, boxrule=0.8pt, arc=2mm, left=5pt, right=5pt
  }
}

% ─── Encabezados ───
\pagestyle{fancy}
\fancyhf{}
\fancyhead[C]{\small\color{gray} \guiaTitulo{} --- \guiaCurso{}}
\fancyfoot[C]{\small\color{gray} Página \thepage}
\renewcommand{\headrulewidth}{0.4pt}

% ─── Títulos ───
\newcommand{\tituloseccion}[1]{%
  \noindent\colorbox{azulprincipal}{%
    \makebox[\dimexpr\textwidth-2\fboxsep\relax][l]{%
      \color{white}\normalfont\Large\bfseries\ SECCIÓN \thesection: #1%
    }%
  }%
}
\titleformat{\section}
  {}{}{0em}
  {\tituloseccion}
\titleformat{\subsection}
  {\normalfont\large\bfseries\color{azulprincipal}}
  {\thesubsection}{0.5em}{}

% ─── Niveles de dificultad ───
\newcommand{\nivelbasico}{%
  \par\smallskip\colorbox{green!60!black}{\color{white}\small\bfseries\ NIVEL BÁSICO\ }\par\smallskip}
\newcommand{\nivelintermedio}{%
  \par\smallskip\colorbox{naranja}{\color{white}\small\bfseries\ NIVEL INTERMEDIO\ }\par\smallskip}
\newcommand{\nivelavanzado}{%
  \par\smallskip\colorbox{rojo}{\color{white}\small\bfseries\ NIVEL AVANZADO\ }\par\smallskip}

% ─── Líneas de respuesta ───
\newcommand{\lineasrespuesta}[1]{%
  \par\vspace{2mm}%
  \foreach \i in {1,...,#1}{\noindent\rule{\textwidth}{0.3pt}\par\vspace{5mm}}%
  \vspace{2mm}%
}

\begin{document}

% ════════════════════════
% PORTADA
% ════════════════════════
\thispagestyle{empty}
\vspace*{3cm}
\begin{center}
  {\Huge\bfseries\color{azulprincipal} GUÍA DE TRABAJO}\\[8pt]
  {\LARGE\bfseries\color{azuloscuro} \guiaTitulo}\\[15pt]
  {\color{azulprincipal}\rule{8cm}{2pt}}\\[15pt]
  {\Large \guiaUnidad{} \textbar{} \guiaCurso}\\[8pt]
  {\large Alumno/a: \guiaAlumno}\\[4pt]
  {\large \guiaSemana{} --- \guiaSemestre}\\[25pt]
  {\itshape\color{gray}
  \guiaSubtitulo}
\end{center}

\vspace{2cm}
\noindent\textbf{Contenidos:}
\begin{enumerate}[leftmargin=2cm]
  \item Motivación: la pregunta inversa de una potencia.
  \item Definición de logaritmo y sus restricciones.
  \item Cálculo de logaritmos simples (bases 2, 3 y 10).
  \item Casos especiales deducidos desde la definición.
  \item Ejercicios por nivel (básico, intermedio, avanzado).
\end{enumerate}
\newpage

% ════════════════════════
% SECCIÓN 1
% ════════════════════════
\section{MOTIVACIÓN: LA PREGUNTA INVERSA DE UNA POTENCIA}

\begin{tcolorbox}[recuerda]
\textbf{Potencia:} repetir una multiplicación.
\[
  2^5 = 2 \cdot 2 \cdot 2 \cdot 2 \cdot 2 = 32
\]
En $2^5$ llamamos a $2$ la \textbf{base} y a $5$ el \textbf{exponente}. El resultado ($32$) es el \textbf{valor de la potencia}.
\end{tcolorbox}

Si me preguntan \emph{``¿cuánto vale $2^5$?''}, respondo $32$. Eso ya lo sabes.\\[4pt]
Pero ahora piensa al revés: \emph{``¿a qué exponente hay que elevar $2$ para obtener $32$?''}. La respuesta sigue siendo $5$, pero la \textbf{pregunta} cambió por completo. Esta nueva pregunta es la que responde el \textbf{logaritmo}.

\begin{tcolorbox}[ejemplo, title=Tres preguntas del mismo tipo]
\begin{itemize}[nosep]
  \item ¿A qué exponente elevo $2$ para obtener $8$? $\longrightarrow$ a $3$ (porque $2^3 = 8$).
  \item ¿A qué exponente elevo $3$ para obtener $81$? $\longrightarrow$ a $4$ (porque $3^4 = 81$).
  \item ¿A qué exponente elevo $10$ para obtener $1000$? $\longrightarrow$ a $3$ (porque $10^3 = 1000$).
\end{itemize}
\end{tcolorbox}

\noindent\textbf{Idea central:} el logaritmo \emph{despeja el exponente}.

% ════════════════════════
% SECCIÓN 2
% ════════════════════════
\section{DEFINICIÓN DE LOGARITMO}

\begin{tcolorbox}[definicion, title=Definición]
Sean $b > 0$, $b \ne 1$ y $a > 0$. El \textbf{logaritmo en base $b$ de $a$}, denotado $\log_b(a)$, es el \textbf{exponente} al que hay que elevar la base $b$ para obtener el argumento $a$.
\[
  \boxed{\; \log_b(a) = c \;\iff\; b^{c} = a \;}
\]
Esta doble flecha significa que las dos expresiones dicen lo mismo: son dos formas de escribir la misma relación entre tres números.
\end{tcolorbox}

\subsection{La traducción entre las dos formas}

Toda igualdad con un logaritmo puede reescribirse como una igualdad con una potencia, y al revés. Esto es lo más importante de la semana.

\begin{tcolorbox}[formula]
\[
  \log_b(a) = c \;\;\Longleftrightarrow\;\; b^{c} = a
\]
\[
  \underbrace{b}_{\text{base}} \quad \underbrace{a}_{\text{argumento}} \quad \underbrace{c}_{\text{resultado del logaritmo } = \text{ exponente}}
\]
\end{tcolorbox}

\begin{tcolorbox}[ejemplo, title=Traducir en ambas direcciones]
\begin{itemize}[nosep]
  \item $\log_2(8) = 3 \;\iff\; 2^{3} = 8$
  \item $\log_3(81) = 4 \;\iff\; 3^{4} = 81$
  \item $\log_{10}(1000) = 3 \;\iff\; 10^{3} = 1000$
  \item $\log_5(25) = 2 \;\iff\; 5^{2} = 25$
\end{itemize}
\end{tcolorbox}

\subsection{¿Por qué las restricciones?}

La definición exige $b > 0$, $b \ne 1$ y $a > 0$. No son caprichos: cada una se rompe si la quitamos.

\begin{tcolorbox}[propiedad, title={Por qué $b > 0$ y $b \ne 1$}]
\begin{itemize}[nosep]
  \item Si $b < 0$: las potencias de base negativa cambian de signo con el exponente ($(-2)^1 = -2$, $(-2)^2 = 4$, $(-2)^3 = -8$, \ldots). No hay forma de que el logaritmo devuelva un resultado único y ordenado.
  \item Si $b = 0$: $0^c = 0$ para todo $c > 0$, así que $\log_0(a)$ no tiene sentido.
  \item Si $b = 1$: $1^c = 1$ para \textbf{cualquier} $c$. Entonces $\log_1(1)$ tendría infinitas respuestas, y $\log_1(5)$ ninguna.
\end{itemize}
\end{tcolorbox}

\begin{tcolorbox}[propiedad, title={Por qué $a > 0$}]
Porque si $b > 0$, entonces $b^c$ siempre es positivo, sin importar el signo de $c$. Nunca podrá ser cero ni negativo. Por lo tanto \textbf{el argumento de un logaritmo debe ser positivo}: expresiones como $\log_2(0)$ o $\log_2(-4)$ \textbf{no están definidas}.
\end{tcolorbox}

\newpage

% ════════════════════════
% SECCIÓN 3
% ════════════════════════
\section{CÁLCULO DE LOGARITMOS SIMPLES}

Cuando el argumento es una potencia exacta de la base, el logaritmo se calcula \textbf{mentalmente}. La estrategia es siempre la misma: preguntarse \emph{``¿a qué exponente elevo la base para llegar al argumento?''}.

\begin{tcolorbox}[ejemplo, title={Ejemplo paso a paso: $\log_2(16)$}]
\begin{enumerate}[nosep]
  \item Pregunta: ¿a qué exponente elevo $2$ para obtener $16$?
  \item Pruebo: $2^1 = 2$, $2^2 = 4$, $2^3 = 8$, $2^{4} = 16$. \checkmark
  \item Respuesta: $\log_2(16) = 4$.
\end{enumerate}
\end{tcolorbox}

\subsection{Tabla de cálculos inmediatos}

\begin{center}
\begin{tabular}{@{}cccc@{}}
\toprule
\textbf{Logaritmo} & \textbf{Pregunta} & \textbf{Potencia equivalente} & \textbf{Resultado} \\
\midrule
$\log_2(8)$    & ¿$2$ a qué da $8$?     & $2^{3} = 8$      & $3$ \\
$\log_2(32)$   & ¿$2$ a qué da $32$?    & $2^{5} = 32$     & $5$ \\
$\log_3(9)$    & ¿$3$ a qué da $9$?     & $3^{2} = 9$      & $2$ \\
$\log_3(27)$   & ¿$3$ a qué da $27$?    & $3^{3} = 27$     & $3$ \\
$\log_{10}(100)$ & ¿$10$ a qué da $100$? & $10^{2} = 100$  & $2$ \\
$\log_{10}(10000)$ & ¿$10$ a qué da $10\,000$? & $10^{4} = 10\,000$ & $4$ \\
\bottomrule
\end{tabular}
\end{center}

% ════════════════════════
% SECCIÓN 4
% ════════════════════════
\section{TRES CASOS ESPECIALES}

Estos tres casos \textbf{no se memorizan}: se \textbf{deducen} desde la definición. Al final de la semana deberías ser capaz de explicarlos sin mirar.

\begin{tcolorbox}[propiedad, title={Caso 1: $\log_b(1) = 0$}]
\textbf{Pregunta:} ¿a qué exponente elevo $b$ para obtener $1$?\\
\textbf{Respuesta:} a $0$, porque \emph{cualquier} base no nula elevada a $0$ vale $1$.
\[
  b^{0} = 1 \;\Longrightarrow\; \log_b(1) = 0
\]
\end{tcolorbox}

\begin{tcolorbox}[propiedad, title={Caso 2: $\log_b(b) = 1$}]
\textbf{Pregunta:} ¿a qué exponente elevo $b$ para obtener $b$?\\
\textbf{Respuesta:} a $1$, porque $b^{1} = b$.
\[
  b^{1} = b \;\Longrightarrow\; \log_b(b) = 1
\]
\end{tcolorbox}

\begin{tcolorbox}[propiedad, title={Caso 3: $\log_b(b^{n}) = n$}]
\textbf{Pregunta:} ¿a qué exponente elevo $b$ para obtener $b^{n}$?\\
\textbf{Respuesta:} a $n$, porque la respuesta está \emph{literalmente escrita} en el argumento.
\[
  b^{n} = b^{n} \;\Longrightarrow\; \log_b(b^{n}) = n
\]
\textit{Este es el caso más potente: todo logaritmo ``se deja ver'' cuando el argumento ya está escrito como potencia de la base.}
\end{tcolorbox}

\newpage

% ════════════════════════
% SECCIÓN 5 — EJERCICIOS
% ════════════════════════
\section{EJERCICIOS}

\nivelbasico

\textbf{(1)} Traduce cada logaritmo a su forma de potencia equivalente (no hace falta calcular el resultado todavía).
\begin{enumerate}[label=\alph*),leftmargin=1cm]
  \item $\log_2(16) = c$ \hfill $\Longleftrightarrow$ \hfill \rule{4cm}{0.4pt}
  \item $\log_3(27) = c$ \hfill $\Longleftrightarrow$ \hfill \rule{4cm}{0.4pt}
  \item $\log_{10}(100) = c$ \hfill $\Longleftrightarrow$ \hfill \rule{4cm}{0.4pt}
  \item $\log_5(125) = c$ \hfill $\Longleftrightarrow$ \hfill \rule{4cm}{0.4pt}
\end{enumerate}

\textbf{(2)} Calcula mentalmente.
\begin{enumerate}[label=\alph*),leftmargin=1cm]
  \item $\log_2(8) = $ \rule{2cm}{0.4pt}
  \item $\log_2(64) = $ \rule{2cm}{0.4pt}
  \item $\log_3(9) = $ \rule{2cm}{0.4pt}
  \item $\log_3(81) = $ \rule{2cm}{0.4pt}
  \item $\log_{10}(10) = $ \rule{2cm}{0.4pt}
  \item $\log_{10}(1000) = $ \rule{2cm}{0.4pt}
  \item $\log_5(25) = $ \rule{2cm}{0.4pt}
  \item $\log_4(16) = $ \rule{2cm}{0.4pt}
\end{enumerate}

\textbf{(3)} Aplica los casos especiales.
\begin{enumerate}[label=\alph*),leftmargin=1cm]
  \item $\log_7(1) = $ \rule{2cm}{0.4pt}
  \item $\log_{12}(12) = $ \rule{2cm}{0.4pt}
  \item $\log_5(5^{4}) = $ \rule{2cm}{0.4pt}
  \item $\log_2(2^{10}) = $ \rule{2cm}{0.4pt}
\end{enumerate}

\nivelintermedio

\textbf{(4)} Encuentra el valor de $x$ (pista: traduce a potencia).
\begin{enumerate}[label=\alph*),leftmargin=1cm]
  \item $\log_3(x) = 4 \;\Longrightarrow\; x = $ \rule{2cm}{0.4pt}
  \item $\log_2(x) = 6 \;\Longrightarrow\; x = $ \rule{2cm}{0.4pt}
  \item $\log_x(49) = 2 \;\Longrightarrow\; x = $ \rule{2cm}{0.4pt}
  \item $\log_x(1000) = 3 \;\Longrightarrow\; x = $ \rule{2cm}{0.4pt}
\end{enumerate}

\textbf{(5)} Calcula. No son inmediatos: ten que darte cuenta de que el argumento \emph{es} una potencia de la base, aunque no se vea a primera vista.
\begin{enumerate}[label=\alph*),leftmargin=1cm]
  \item $\log_2\!\left(\tfrac{1}{8}\right)$. \hfill \textit{Pista: $\tfrac{1}{8} = 2^{-3}$.}
  \lineasrespuesta{2}
  \item $\log_3\!\left(\tfrac{1}{81}\right)$.
  \lineasrespuesta{2}
  \item $\log_2(\sqrt{2})$. \hfill \textit{Pista: $\sqrt{2} = 2^{1/2}$.}
  \lineasrespuesta{2}
\end{enumerate}

\nivelavanzado

\textbf{(6)} Decide si cada expresión está definida. Si \textbf{no} lo está, explica por qué en una línea.
\begin{enumerate}[label=\alph*),leftmargin=1cm]
  \item $\log_2(-4)$
  \lineasrespuesta{2}
  \item $\log_1(5)$
  \lineasrespuesta{2}
  \item $\log_{10}(0)$
  \lineasrespuesta{2}
  \item $\log_{-2}(4)$
  \lineasrespuesta{2}
\end{enumerate}

\textbf{(7)} Justifica en tus palabras (2--3 líneas) por qué $\log_b(b^{n}) = n$ para cualquier exponente $n$, incluyendo el caso $n = 0$ y $n$ negativo.
\lineasrespuesta{4}

\newpage

% ════════════════════════
% RESPUESTAS SELECCIONADAS
% ════════════════════════
\section*{Respuestas seleccionadas}
\addcontentsline{toc}{section}{Respuestas seleccionadas}

\begin{tcolorbox}[propiedad, title=Soluciones (básico e intermedio)]
\begin{itemize}[nosep]
  \item \textbf{(1)} a) $2^{c} = 16$; \; b) $3^{c} = 27$; \; c) $10^{c} = 100$; \; d) $5^{c} = 125$.
  \item \textbf{(2)} a) $3$; \; b) $6$; \; c) $2$; \; d) $4$; \; e) $1$; \; f) $3$; \; g) $2$; \; h) $2$.
  \item \textbf{(3)} a) $0$; \; b) $1$; \; c) $4$; \; d) $10$.
  \item \textbf{(4)} a) $x = 81$; \; b) $x = 64$; \; c) $x = 7$; \; d) $x = 10$.
  \item \textbf{(5)} a) $-3$; \; b) $-4$; \; c) $\tfrac{1}{2}$.
\end{itemize}
\end{tcolorbox}

\begin{tcolorbox}[recuerda]
Las respuestas del \textbf{nivel avanzado} (6 y 7) se revisan en clase: la idea es que escribas tu propio razonamiento, no que copies una respuesta tipo.
\end{tcolorbox}

\end{document}
